\clearpage
\section*{Question 2}
\addcontentsline{toc}{section}{Question 2}

\textbf{Problem Statement:} Suppose you have an iterative algorithm for finding a fixed point of a function $F(x)$ on a certain interval. Write a short text (2-4 paragraphs) explaining the following:

\begin{enumerate}
\item How do you decide when to stop iterating?
\item What is your strategy for making sure that the computed solution is close to the exact solution?
\end{enumerate}

DO NOT use the standard methodology and estimates from the textbook. Try to find other ways of accomplishing the goals above. Discuss briefly the advantages and disadvantages of your error control approach.

\vspace{1em}

\textbf{Solution:}

\subsection*{Assumptions:}

\begin{enumerate}
\item $F(x)$ is a contractive function in the given interval (I have the expression and can use it to generate $x_n$ as $n = 1, 2, 3, \ldots$ and some $x_0$).
\item I start with an $x_0$ `close enough' to `r', so that I can assume that $F(x)$ will converge to $r$ with enough iterations, either because it was in its `region of contraction' or even otherwise. Else by running my solvers/algorithms repeatedly I can switch $x_0$ till I get a reliable convergence to some fixed point $r$.
\end{enumerate}

\subsection*{Approach:}

From knowledge of $F(x)$, I can determine the order of convergence via Taylor's formula.

\subsection*{Step 1: Generate Table of Fixed Point Iterations}

Since I don't know the exact fixed point $r$, I cannot directly compute the errors $e_n = |x_n - r|$. However, I can see whether my consecutive errors are following (approximately) order of convergence $p$.

\subsection*{Step 2: Table Generation}

Generate a table of $n$ vs $x_n$ values by recursively using $x_{n+1} = F(x_n)$:

\begin{center}
\begin{tabular}{|c|c|}
\hline
$n$ & $x_n$ \\
\hline
0 & $x_0$ \\
1 & $x_1$ \\
2 & $x_2$ \\
$\vdots$ & $\vdots$ \\
$n$ & $x_n$ \\
$\vdots$ & $\vdots$ \\
\hline
\end{tabular}
\end{center}

\subsection*{Step 3: Estimating the Order of Convergence $p$}

For a convergent sequence with order of convergence $p$, I have:
$$\lim_{n \to \infty} \frac{|x_{n+1} - r|}{|x_n - r|^p} = C$$

where $C$ is a constant and $r$ is the unknown fixed point.

Since I don't know $r$, I can use the approximation that for some sufficiently large $n$, $x_n \approx r$. This gives me:

$$\frac{|x_{n+2} - x_{n+1}|}{|x_{n+1} - x_n|^p} \approx C$$

Taking logarithms of consecutive ratios:
$$\ln\left(\frac{|x_{n+2} - x_{n+1}|}{|x_{n+1} - x_n|}\right) \approx p \ln\left(\frac{|x_{n+1} - x_n|}{|x_n - x_{n-1}|}\right)$$

Rearranging to solve for $p$:
\begin{align}
p &\approx \frac{\ln\left(\frac{|x_{n+2} - x_{n+1}|}{|x_{n+1} - x_n|}\right)}{\ln\left(\frac{|x_{n+1} - x_n|}{|x_n - x_{n-1}|}\right)} \label{eq:p-order-approximate}
\end{align}
% $$p \approx \frac{\ln\left(\frac{|x_{n+2} - x_{n+1}|}{|x_{n+1} - x_n|}\right)}{\ln\left(\frac{|x_{n+1} - x_n|}{|x_n - x_{n-1}|}\right)} \label{eq:p-order-approximate}$$

This formula allows me to estimate $p$ using only the computed values $x_{n-1}$, $x_n$, $x_{n+1}$, and $x_{n+2}$ without knowing the exact fixed point $r$.

\textbf{When do I know that the computed solution is close enough to the exact solution?}\\
Once the RHS of Equation \ref{eq:p-order-approximate} stabilizes to a constant value (i.e., the estimated $p$ does not change significantly with further iterations), I can be confident that I'm `heading towards' the fixed point $r$ i.e. my FPI iterations are in the `region of contraction'.

\textbf{How do you decide when to stop iterating?}\\
Once `sufficient closeness' to $r$ is determined, I can employ a tolerance on the absolute (approximate) error value, say $\epsilon$,. I can then stop the iterations when my approximate $e_n$ falls below this threshold and conclude that my approximate $x_n \approx r$. $\quad \square$